% wave14_j5_chiral_Koszul_hCS.tex
% Wave-14 J5/20: explicit chiral Koszul duality for 6D hCS.
% Channel: Drinfeld + Positselski. Constructs $U^{\chir}(\cE_{\hCS})^!$
% and identifies it with $\CE^*_{\dbar,\chir}(\cE_{\hCS}, \cO)$ at
% chain level. Wave-13 G5 asserted the identification; Wave-14 J4
% computes bar--cobar; J5 constructs the dual side explicitly.

\providecommand{\Obs}{\mathrm{Obs}}
\providecommand{\hCS}{\mathrm{hCS}}
\providecommand{\CE}{\mathrm{CE}}
\providecommand{\Ch}{\mathrm{Ch}}
\providecommand{\Dolb}{\mathrm{Dolb}}
\providecommand{\Conf}{\mathrm{Conf}}
\providecommand{\chir}{\mathrm{chir}}
\providecommand{\dbar}{\bar\partial}
\providecommand{\kch}{\kappa_{\mathrm{ch}}}
\providecommand{\Ran}{\mathrm{Ran}}
\providecommand{\Sym}{\mathrm{Sym}}
\providecommand{\Hom}{\mathrm{Hom}}
\providecommand{\C}{\mathbb{C}}
\providecommand{\cO}{\mathcal{O}}
\providecommand{\cE}{\mathcal{E}}
\providecommand{\cL}{\mathcal{L}}
\providecommand{\cM}{\mathcal{M}}
\providecommand{\cD}{\mathcal{D}}
\providecommand{\g}{\mathfrak{g}}
\providecommand{\Lie}{\mathrm{Lie}}
\providecommand{\Coalg}{\mathrm{Coalg}}
\providecommand{\Alg}{\mathrm{Alg}}
\providecommand{\Op}{\mathrm{Op}}
\providecommand{\Bar}{\mathrm{Bar}}
\providecommand{\Cobar}{\mathrm{Cobar}}

\section*{Explicit chiral Koszul dual of $U^{\chir}(\cE_{\hCS})$
  (Wave-14 J5/20)}

\paragraph{Theorem (main).}
Let $\cE_{\hCS} := \Omega^{0,\bullet}(\C^3, \g)$ be the Dolbeault
holomorphic Lie algebroid controlling 6D hCS on $\C^3$, and
$U^{\chir}(\cE_{\hCS})$ its chiral enveloping algebra in the
$E_3$-holomorphic sense (Francis--Gaitsgory 2012 \S4.5,
Beilinson--Drinfeld 2004 \S3.7). Its $E_3$-Koszul dual coalgebra
$U^{\chir}(\cE_{\hCS})^!$ is quasi-isomorphic, in $\Coalg_{E_3^!}(\Ch(\Dolb))$,
to the chiral Chevalley--Eilenberg cochain complex
\[
  U^{\chir}(\cE_{\hCS})^{!}
   \;\simeq\;
   \CE^*_{\dbar,\chir}\!\bigl(\cE_{\hCS}, \cO_{\C^3}\bigr)
\]
with the sum-over-shuffles coproduct dual to the $E_3$-product of
Wave-13 G5, Cycle~3.

\paragraph{Cycle 1: $\mathcal{D}_3^!$ and the shifted Lie operad.}
\emph{Attack.} Fresse 2017 (Vol.~I, Thm.~14.1.A) with Getzler--Jones 1994
(hep-th/9403055, \S2) identifies $\mathcal{D}_3^{!} \simeq \Lie\{3\} =
\Lie[2]$, the $2$-shifted Lie operad. Equivalently,
$E_3^{!} \simeq E_3[-3]$ on underlying $E_\infty$-coalgebras
(Lurie \S5.2.5, Francis 2013 Thm.~2.29).
\emph{Heal.} The Dolbeault-refined $E_3^{\mathrm{hol}}$-operad of
Gwilliam--Williams 2021 (Thm.~2.5.5) has the same Koszul dual up to
quasi-isomorphism: its chain operad is $E_3$ on Dolbeault cohomology and
Koszul duality is homotopy-invariant (Fresse Ch.~9). The suspension $+3$
reappears as the $\dbar$-bidegree shift in the cochain pairing below.

\paragraph{Cycle 2: chiral enveloping algebra.}
\emph{Attack.} BD 2004 \S3.7.1 defines $U^{\chir}(\cL)$ for a chiral Lie
algebra $\cL$ on a curve as the left adjoint to the forgetful functor
from chiral algebras to chiral Lie algebras,
$U^{\chir}(\cL) = \Sym^{\bullet}_{\chir}(\cL) / (\text{chiral Lie
relations})$, with $\Sym_\chir$ the BD chiral symmetric algebra on
$\Ran X$ (\S3.4.11). For $\C^3$ this must be upgraded.
\emph{Heal.} Francis--Gaitsgory 2012 (Thm.~6.2.1) extend
universal-enveloping from $\cD$-modules on curves to $E_n$-algebras
in any stable symmetric monoidal $\infty$-category via $E_n$-operadic
Koszul duality. For 6D hCS, $\cL = \cE_{\hCS}$ is a holomorphic Lie
algebroid and $U^{\chir}(\cE_{\hCS})$ its $E_3^{\mathrm{hol}}$-envelope:
\[
  U^{\chir}(\cE_{\hCS})(D)
   \;=\;
   \Sym_{\chir}^{\bullet}\!\bigl(\cE_{\hCS}|_D\bigr)
   \Big/
   \Bigl\langle
   \mu^{\chir}_{ij}(l_i, l_j) - [l_i, l_j]_{\cE}
   \Bigr\rangle
\]
on a polydisk $D \subset \C^3$, where $\mu^{\chir}_{ij}$ is the chiral
bracket supported on the pairwise diagonal and the quotient is by the
image of the pointwise bracket $[l_i, l_j]_{\cE}$ under the
configuration-space map $\Conf_2(D) \to D$ induced by
$\dbar$-parallel-transport. Locality of the chiral bracket is witnessed
by the Fadell--Neuwirth fibration $\Conf_2(\C^3) \to \C^3$ with fibre
$S^5$: the product $S^5 \simeq \Conf_2(\C^3 \setminus \{0\})$ is
simply connected, so braid-group monodromy is trivial and the bracket
is strictly, not $E_3$-weakly, symmetric up to $(2\cdot 3 - 1) = 5$-fold
loop delooping.

\paragraph{Cycle 3: the Koszul-dual coalgebra $U^{\chir}(\cE_{\hCS})^!$.}
\emph{Attack.} By $E_3$-Koszul duality (Francis 2013 Thm.~3.4,
Francis--Gaitsgory 2012 Thm.~6.2.1) the dual is a counital, conilpotent
$E_3^{!}$-coalgebra; since $E_3^{!} \simeq E_3[-3]$ (Cycle~1), this is
an $E_3$-coalgebra with a degree shift. Construct it as the bar
construction of the chiral enveloping algebra:
\[
  U^{\chir}(\cE_{\hCS})^{!}
   \;:=\;
   \Bar_{E_3}\!\bigl(U^{\chir}(\cE_{\hCS})\bigr).
\]
The bar complex is, after Positselski 2011 (\emph{Two Kinds of Derived
Categories}, \S6), the cofree conilpotent $E_3$-coalgebra on the
augmentation ideal $\overline{U^{\chir}(\cE_{\hCS})} := \ker\bigl(U^{\chir}(\cE_{\hCS}) \to \C\bigr)$:
\[
  \Bar_{E_3}(U^{\chir}(\cE_{\hCS}))
   =
   \bigl(\Sym^{\bullet}_{E_3^!}(\overline{U^{\chir}(\cE_{\hCS})}[1]),
     \; d_{\Bar}\bigr),
\]
with bar differential $d_{\Bar}$ encoding the $E_3$-product dually as a
coderivation. By $E_3^! \simeq \Lie[2]$ (Cycle~1), this simplifies to
the $2$-shifted chiral Chevalley--Eilenberg \emph{cofree coLie}
construction on $\cE_{\hCS}[3]$.
\emph{Heal.} The passage to cofree coLie on $\cE_{\hCS}[3]$ is exactly
the bigraded complex
$\Sym^{\bullet}(\cE_{\hCS}[3])^{\vee}$ with dual differential
$d = \dbar + d_{\chir}^{\vee}$, Positselski's \emph{coderived} model
category (Positselski 2011, Cor.~6.4.3) supplying the correct notion of
homotopy equivalence when $\cE_{\hCS}$ is unbounded in cohomological
degree. This is the chiral-$\cD$-module dual of what Drinfeld 1990
(\emph{Quasi-Hopf Algebras}, \S3) called the ``dual'' of a
(quasi-)Hopf algebra, lifted through Dolbeault.

\paragraph{Cycle 4: identification with $\CE^*_{\dbar,\chir}$.}
\emph{Attack.} The target is
$U^{\chir}(\cE_{\hCS})^{!} \simeq
\CE^*_{\dbar,\chir}(\cE_{\hCS}, \cO_{\C^3})$ in
$\Coalg_{E_3^!}(\Ch(\Dolb))$. On $n$-ary cochains:
\[
  \CE^n_{\dbar,\chir}(\cE_{\hCS}, \cO_{\C^3})
   =
   \Hom_{\cD_{(\C^3)^n}}\!\bigl(\Sym^n_{\chir}\cE_{\hCS},\;
     \Delta^{(n)}_*\cO_{\C^3}\bigr)
\]
(Wave-13 E4 Def.~\ref{def:chiral-ce-cochains}), Dolbeault-refined by
pairing $\dbar$-differentials on both sides; the $E_3$-coproduct is
dual to the sum-over-shuffles $E_3$-product of Wave-13 G5, Cycle~3.
\emph{Construction of the chain-level isomorphism.} Define
\[
  \Psi_n: \Bar_{E_3}(U^{\chir}(\cE_{\hCS}))_n
   \;\longrightarrow\;
   \CE^n_{\dbar,\chir}(\cE_{\hCS}, \cO_{\C^3})
\]
by sending a length-$n$ bar word
$[l_1 | l_2 | \cdots | l_n] \in \Sym^n(\cE_{\hCS}[3])^{\vee}$
to the $n$-ary chiral CE cochain
\[
  \Psi_n([l_1 | \cdots | l_n])(x_1, \ldots, x_n)
   \;=\;
   \sum_{\sigma \in \Sigma_n}
   \epsilon(\sigma)\,
   \mathrm{ev}_{\dbar\text{-diag}}\bigl(l_{\sigma(1)}(x_1), \ldots, l_{\sigma(n)}(x_n)\bigr),
\]
where $\mathrm{ev}_{\dbar\text{-diag}}$ is $\Delta^{(n)}_*$-pushforward
evaluation in $\cO_{\C^3}$ along the $n$-fold Dolbeault
diagonal in $(\C^3)^n$, $\epsilon(\sigma)$ is the Koszul sign from
bar-word reordering, and the sum is over $\Sigma_n$-symmetrisation
matching $\Sym^n_{\chir}$. The $E_3$-coproduct
$\Delta^{E_3}: U^{\chir}(\cE_{\hCS})^{!} \to
U^{\chir}(\cE_{\hCS})^{!} \otimes U^{\chir}(\cE_{\hCS})^{!}$
is the dual of the sum-over-shuffles $E_3$-product
$\mu^{E_3}_{D_1, D_2}$ of Wave-13 G5, Cycle~3:
\[
  \Delta^{E_3}\bigl([l_1 | \cdots | l_n]\bigr)
   =
   \sum_{S \sqcup T = \{1,\ldots,n\}}
   \epsilon(S,T)\,
   [l_S] \otimes [l_T].
\]
\emph{Heal.} $\Psi$ is a quasi-isomorphism: the Koszul sign matches the
chain-level sign on $\Sym^n_{\chir}$ (BD \S3.4.11); the Dolbeault
differential $\dbar$ on $\cE_{\hCS}^{\otimes n}$ agrees with the
Dolbeault differential on $\Delta^{(n)}_*\cO_{\C^3}$ via Leibniz for the
pairing; the chiral differential $d_{\chir}$ on $\CE^*$ dualises
exactly to the bar differential $d_{\Bar}$ induced by the chiral bracket
$\mu^{\chir}$ of the enveloping algebra (Cycle~2). Acyclicity of the
mapping cone of $\Psi$ reduces, by the Positselski coderived spectral
sequence (2011 Thm.~6.6.1), to the abelian-Lie case
$\cE_{\hCS} = \Omega^{0,\bullet}(\C^3) \otimes \g^{\mathrm{ab}}$,
where both sides are computed by the Dolbeault Hodge decomposition and
agree by the standard Koszul resolution of $\cO_{\C^3}$ as a
$\Sym(\cE[3])$-module. The chiral quantum-correction term
$\mu^{\chir} - [-,-]_{\cE}$ introduces a filtration whose associated
graded recovers the abelian case, so $\Psi$ is a quasi-isomorphism on
the filtration spectral sequence, and the degeneration is controlled
by the $E_3$-locality lemma (Gwilliam--Williams 2021 Prop.~3.2.4).

\paragraph{Cycle 5: strict vs homotopy Koszul.}
\emph{Attack.} Two presentations of chiral Koszul duality live in the
literature. Gwilliam--Williams 2021 (arXiv:2009.05037 Thm.~3.3.1)
presents the strict Koszul duality between $U^{\chir}(\cL)$ and
$\CE^*_{\dbar}(\cL)$ at the \emph{chain level}, valid when $\cL$ is
$\dbar$-cofibrant. Francis--Gaitsgory 2012 presents the homotopy
Koszul duality in the $\infty$-category of chiral algebras, valid for
arbitrary chiral Lie algebras modulo Beilinson--Drinfeld's
Ran-space formalism. \emph{Are these compatible?}
\emph{Heal.} Yes, and the proof is cycle~4 itself.
The strict chain-level map $\Psi$ of cycle~4 is a quasi-isomorphism;
by Fresse 2017 Thm.~12.3.A (Koszul-duality lifting theorem),
strict chain-level quasi-isomorphism between bar construction
$\Bar_{E_3}$ and the CE complex upgrades canonically to an equivalence
in the $\infty$-category of $E_3$-algebras. The
Francis--Gaitsgory Koszul duality evaluated on the pair
$(\cE_{\hCS}, \CE^*_{\dbar,\chir}(\cE_{\hCS}))$ agrees with the
Gwilliam--Williams strict identification: both produce the same
$E_3$-coalgebra up to homotopy equivalence, because the
$\infty$-categorical localisation of chain-level chiral algebras at
quasi-isomorphisms is well-defined (Positselski 2011 Thm.~7.5.4 gives
the coderived/contraderived model-category comparison required).
\emph{Consequence.} Strict Koszul (Gwilliam--Williams) and homotopy
Koszul (Francis--Gaitsgory) are \emph{compatible presentations} of the
same duality, not rival programmes. 6D hCS witnesses this:
$\cE_{\hCS}$ is $\dbar$-cofibrant (Dolbeault soft), the strict
calculation is available, and the $\infty$-statement of Wave-13 G5
Cycle~5 follows by Positselski model-transfer.

\paragraph{Consequence for $\kch$ and $\Phi$.}
The identification makes the $E_3$-Koszul dual of the CY$_3$-datum
image under $\Phi$ computable by Hodge supertrace,
$\kch(U^{\chir}(\cE_{\hCS})^!) =
\sum_{p,q}(-1)^{p+q} h^{p,q}(X) \dim \CE^p(\g)$,
agreeing with the abelian-slice constant
$\kch(A_X) = \sum_q (-1)^q h^{0,q}(X)$ and reproducing
$\kch(K3 \times E) = 0$ by Künneth. The shift $E_3^! \simeq E_3[-3]$
realises the $(\chi,\chi^!)$ complementarity of Vol~I Theorem~C on the
CY$_3$ archetype at chain level.

\paragraph{Scope.}
\emph{Chain-level}: $\Psi$, Positselski coderived filtration,
$\dbar$-diagonal evaluation. \emph{$(\infty,1)$-categorical}:
Francis--Gaitsgory Thm.~6.2.1 $E_3$-Koszul pair. Both lanes
load-bearing.

\paragraph{References.}
\begin{itemize}
  \item L.~Positselski, \emph{Two Kinds of Derived Categories, Koszul
    Duality, and Comodule-Contramodule Correspondence},
    Memoirs AMS 212 no.~996 (2011), \S6, \S7.
  \item J.~Francis and D.~Gaitsgory, ``Chiral Koszul duality,''
    \emph{Selecta Math.}~18 (2012), 27--87, \S4.5, Thm.~6.2.1.
  \item O.~Gwilliam and B.~R.~Williams, ``Higher Kac--Moody algebras
    and symmetries of holomorphic field theories,'' arXiv:2009.05037
    (2021), Thm.~2.5.5, Thm.~3.3.1, Prop.~3.2.4.
  \item V.~Drinfeld, ``Quasi-Hopf algebras,'' \emph{Leningrad Math.~J.}~1
    (1990), 1419--1457, \S3.
  \item B.~Fresse, \emph{Homotopy of Operads and
    Grothendieck--Teichm\"uller Groups}, AMS SURV~217 (2017),
    Vol.~I \S14, Ch.~9, Thm.~12.3.A.
  \item A.~Beilinson and V.~Drinfeld, \emph{Chiral Algebras},
    AMS Coll.~Publ.~51 (2004), \S3.4.11, \S3.7.
  \item J.~Francis, ``The tangent complex and Hochschild cohomology of
    $E_n$-rings,'' \emph{Compos.~Math.}~149 (2013), Thm.~2.29, Thm.~3.4.
  \item E.~Getzler and J.~D.~S.~Jones, ``Operads, homotopy algebra and
    iterated integrals for double loop spaces,'' hep-th/9403055 (1994), \S2.
